\documentclass[11pt,a4paper]{article}

\usepackage[utf8]{inputenc}
\usepackage[T1]{fontenc}
\usepackage[french]{babel}
\usepackage{geometry}
\usepackage{longtable}
\usepackage{array}
\usepackage{booktabs}
\usepackage{hyperref}
\usepackage{xcolor}
\usepackage{listings}
\usepackage{enumitem}
\usepackage{graphicx}
\usepackage{float}

\geometry{margin=2.4cm}
\hypersetup{
  colorlinks=true,
  linkcolor=black,
  urlcolor=blue,
  citecolor=black
}

\lstdefinestyle{trackrcode}{
  basicstyle=\ttfamily\small,
  breaklines=true,
  frame=single,
  columns=fullflexible,
  keepspaces=true,
  showstringspaces=false
}

\setlist[itemize]{noitemsep, topsep=0.25em}
\setlist[enumerate]{noitemsep, topsep=0.25em}
\renewcommand{\arraystretch}{1.25}

\title{\textbf{Rapport technique final -- TrackrAI}}
\author{}
\date{\today}

\begin{document}

\maketitle
\tableofcontents
\newpage

\section{Introduction}

TrackrAI est une application de suivi sportif destinée à centraliser des données issues de séances physiques, à les visualiser et à fournir des éléments d'analyse au sportif ainsi qu'à son encadrement. L'application associe un module matériel à base d'ESP32, une interface web/mobile, une API, une base de données, un serveur de centralisation et un serveur d'analyse vidéo.

Le contexte principal d'utilisation est le suivi d'une séance d'entraînement réalisée par un sportif. Une séance peut être démarrée depuis l'interface applicative, puis alimentée par un module connecté qui transmet des mesures GPS, inertielles, cardiaques et un compteur de pas calculé localement sur l'ESP32 à partir de l'accéléromètre. Les données sont ensuite agrégées, filtrées et transformées en statistiques de séance. Le sportif consulte son historique, ses statistiques et les prévisions cardiaques, compare deux séances et voit son coach référent. Le coach consulte les sportifs qui lui sont affectés, suit une séance active en direct depuis le détail de séance et compare leurs séances avec une lecture orientée progression. L'administrateur gère les utilisateurs, les modules et les affectations coach/sportif.

L'application propose également une analyse vidéo de mouvements sportifs. Le flux vidéo est envoyé à un serveur Python par WebSocket, analysé avec MediaPipe/OpenCV, puis stocké dans un historique intégré à la page d'analyse vidéo. Le squat constitue le mode d'analyse principal. Des modes expérimentaux existent pour le développé couché et le soulevé de terre.

Enfin, un modèle XGBoost prévoit la fréquence cardiaque moyenne de la prochaine séance comparable. La cible est une mesure réellement observée, en bpm, et non un score fabriqué. L'entraînement combine les séances locales de MongoDB et des fenêtres cardiaques et inertielles extraites de l'archive officielle PAMAP2 de l'UCI. Lorsqu'une séance comparable suivante est réalisée, l'application mesure l'erreur entre prévision et observation afin de constituer progressivement un retour propre au sportif.

\section{Structure logicielle et matérielle globale}

\subsection{Vue d'ensemble}

TrackrAI est organisé en plusieurs services indépendants communiquant par HTTP, TCP, WebSocket et MongoDB. La séparation des responsabilités permet de remplacer ou d'étendre une partie de l'application sans réécrire l'ensemble du système.

\begin{figure}[H]
\centering
\begin{verbatim}
ESP32 + GPS + MPU6050 + ceinture cardio BLE
        |
        | TCP : AUTOREGISTER, HELLO, STOREMEASURE
        v
Serveur central Java TCP
        |
        | HTTP interne sécurisé par x-service-secret
        v
API Node.js / Express  <---- HTTP ---- Frontend Vue / Capacitor
        |                                  |
        | Mongoose                         | WebSocket
        v                                  v
MongoDB                         Serveur analyse vidéo Python
                                      |
                                      | TCP : STOREANALYSIS
                                      v
                              Serveur central Java
\end{verbatim}
\caption{Architecture globale de TrackrAI}
\end{figure}

\subsection{Composants principaux}

\begin{longtable}{p{3.2cm}p{3.2cm}p{7.2cm}}
\toprule
\textbf{Composant} & \textbf{Technologie} & \textbf{Rôle} \\
\midrule
Frontend & Vue 3, Vite, Pinia, Chart.js, Leaflet, Capacitor & Interface sportif, coach et administrateur. Il consomme l'API HTTP et envoie les vidéos au serveur d'analyse par WebSocket. \\
API & Node.js, Express, Mongoose & Expose les routes HTTP, applique l'authentification, contrôle les rôles, écrit et lit les données MongoDB, lance les scripts IA. \\
Base de données & MongoDB 6 & Stocke les utilisateurs, modules, chipsets, mesures, séances, statistiques et résultats d'analyse vidéo. \\
Serveur central & Java, sockets TCP, HttpClient & Reçoit les connexions des modules ESP32, gère l'état des sessions côté module et relaie les mesures vers l'API. \\
Serveur d'analyse & Python, websockets, OpenCV, MediaPipe & Reçoit les vidéos, extrait les poses, calcule des métriques de mouvement et transmet le résultat au serveur central. \\
Module matériel & ESP32, GPS, MPU6050, BLE HR & Acquiert les mesures physiques, calcule un compteur de pas embarqué à partir de l'accéléromètre et transmet les données au serveur central pendant une séance active. \\
\bottomrule
\end{longtable}

\section{Contexte logiciel et matériel}

\subsection{Arborescence des sources}

Les sources sont versionnées dans le dépôt Git de l'application. Une reprise du projet se fait en clonant ce dépôt, puis en utilisant les fichiers Docker Compose et les guides de lancement fournis à la racine. Les répertoires principaux sont :

\begin{longtable}{p{4cm}p{10cm}}
\toprule
\textbf{Répertoire} & \textbf{Contenu} \\
\midrule
\texttt{WebTrackr/} & Application frontend Vue 3. Contient les vues, composants, stores Pinia, services HTTP/WebSocket et configuration Vite/Capacitor. \\
\texttt{TrackrAPI/} & API Node.js/Express. Contient les contrôleurs, routes, modèles Mongoose, scripts IA et configuration de démarrage. \\
\texttt{CentralServer/} & Serveur Java TCP et client de test. Contient l'interface \texttt{DataDriver}, les implémentations Mongo/HTTP et la logique de gestion des sockets. \\
\texttt{AnalyzeServer/} & Serveur Python WebSocket d'analyse vidéo. Contient les analyseurs, le stockage temporaire des vidéos et le client TCP vers le serveur central. \\
\texttt{arduino/} & Firmware ESP32 PlatformIO. Contient le sketch de connexion WiFi, acquisition GPS/IMU/BLE et envoi TCP des mesures. \\
\texttt{TrackrAIAcceptanceTests/} & Tests d'acceptation automatisés. \\
\texttt{docs/} & Documentation complémentaire sur l'intégration continue GitLab/Jenkins. \\
\texttt{DEV.md} & Guide de lancement, tests rapides, comptes de démonstration, reset de base, Swagger et procédure ESP32. \\
\bottomrule
\end{longtable}

\subsection{Prérequis de développement}

\begin{longtable}{p{4cm}p{4cm}p{6cm}}
\toprule
\textbf{Outil} & \textbf{Version indicative} & \textbf{Utilisation} \\
\midrule
Docker & 24 ou supérieur & Exécution des services en développement et production. \\
Docker Compose & v2 & Orchestration de MongoDB, API, serveur central, serveur d'analyse et frontend. \\
Node.js & 20 ou supérieur & Exécution de l'API et du frontend lorsqu'ils ne sont pas conteneurisés. \\
MongoDB & 6.0 & Base de données principale. \\
Java & 21 & Compilation et exécution du serveur central. \\
Python & 3.x & Exécution du serveur d'analyse vidéo et scripts IA. \\
PlatformIO & Version récente & Compilation et téléversement du firmware ESP32. \\
\bottomrule
\end{longtable}

\subsection{Dépendances principales}

\subsubsection{Frontend}

Le frontend utilise Vue 3, Vite, Pinia pour l'état applicatif, Vue Router pour la navigation, Axios pour les appels HTTP, Chart.js pour les graphiques et Leaflet pour les cartes. Capacitor fournit l'enveloppe mobile Android. En développement, l'interface est servie par Vite. En production, le build statique est servi par un conteneur Nginx qui joue également le rôle de reverse proxy vers l'API Node et le serveur WebSocket d'analyse vidéo.

\begin{lstlisting}[style=trackrcode]
cd WebTrackr
npm install
npm run dev
npm run build
\end{lstlisting}

\subsubsection{API}

L'API utilise Express, Mongoose, bcryptjs, uuid, validator, dotenv et nodemon en développement. Les scripts IA utilisent Python, pandas, scikit-learn, XGBoost et joblib dans l'image Docker de l'API.

\begin{lstlisting}[style=trackrcode]
cd TrackrAPI
npm install
npm run dev
\end{lstlisting}

\subsubsection{Analyse vidéo}

Le serveur d'analyse vidéo repose sur les bibliothèques Python suivantes :

\begin{lstlisting}[style=trackrcode]
opencv-python
mediapipe==0.10.14
numpy<2
websockets==12.0
\end{lstlisting}

\subsection{Lancement par Docker Compose}

Le lancement recommandé en développement se fait à la racine du dépôt. Cette commande construit les images locales si nécessaire, installe les dépendances dans les conteneurs et démarre l'ensemble des services applicatifs :

\begin{lstlisting}[style=trackrcode]
docker compose -f docker-compose.dev.yml up --build
\end{lstlisting}

Les services exposés en développement sont :

\begin{longtable}{p{4cm}p{3cm}p{7cm}}
\toprule
\textbf{Service} & \textbf{Port} & \textbf{Rôle} \\
\midrule
Frontend Vite & 5173 & Interface web de développement. \\
API Trackr & 4567 & Routes HTTP sous le préfixe \texttt{/trackrapi}. \\
Serveur central & 29000 vers 9000 & Point d'entrée TCP pour les modules ESP32. \\
Serveur analyse & 6000 & WebSocket pour l'upload et l'analyse vidéo. \\
MongoDB & 27017 & Base de données accessible localement. \\
\bottomrule
\end{longtable}

Pour un lancement de production local, le fichier \texttt{docker-compose.yml} construit le frontend puis le sert avec Nginx. L'API, le serveur central, le serveur d'analyse et MongoDB restent isolés dans le réseau Docker :

\begin{lstlisting}[style=trackrcode]
docker compose -f docker-compose.yml up --build -d
\end{lstlisting}

L'arrêt de la pile de production locale se fait par :

\begin{lstlisting}[style=trackrcode]
docker compose -f docker-compose.yml down
\end{lstlisting}

Les variables d'environnement principales sont :

\begin{longtable}{p{5cm}p{9cm}}
\toprule
\textbf{Variable} & \textbf{Description} \\
\midrule
\texttt{MONGODB\_URL} & URL de connexion MongoDB. \\
\texttt{PORT} & Port HTTP de l'API. \\
\texttt{CENTRAL\_HOST}, \texttt{CENTRAL\_PORT} & Adresse du serveur central depuis les services internes. \\
\texttt{SERVICE\_SECRET} & Secret partagé utilisé par les appels internes Java/API. \\
\texttt{JWT\_SECRET} & Secret utilisé par l'API pour signer et vérifier les access tokens JWT. \\
\texttt{USER\_ACCESS\_TOKEN\_TTL\_MINUTES} & Durée de validité de l'access token utilisateur. \\
\texttt{USER\_REFRESH\_TOKEN\_TTL\_DAYS} & Durée de validité du refresh token utilisateur. \\
\texttt{WS\_HOST}, \texttt{WS\_PORT} & Adresse d'écoute du serveur WebSocket d'analyse vidéo. \\
\texttt{MAX\_VIDEO\_BYTES} & Taille maximale d'une vidéo reçue par le serveur d'analyse. \\
\bottomrule
\end{longtable}

\section{Descriptif technique}

\subsection{Frontend Vue / Capacitor}

\subsubsection{Organisation}

Le répertoire \texttt{WebTrackr/src} contient les vues, composants, services et stores. Les services isolent les communications HTTP et WebSocket. Les stores Pinia conservent l'état de session utilisateur, les droits, l'identité du sportif et les données nécessaires aux pages.

\begin{longtable}{p{5cm}p{9cm}}
\toprule
\textbf{Élément} & \textbf{Rôle} \\
\midrule
\texttt{views/} & Pages principales : tableau de bord sportif, historique, détail séance avec mode direct, comparaison de séances, administration, coach, analyse vidéo avec historique intégré. \\
\texttt{components/} & Composants réutilisables : navigation, graphiques, cartes statistiques. \\
\texttt{services/api.service.js} & Client Axios, ajout du JWT dans \texttt{Authorization: Bearer}, compatibilité \texttt{x-session-id}, refresh automatique sur réponse 401. \\
\texttt{services/auth.service.js} & Connexion, déconnexion, refresh token et vérification de session. \\
\texttt{services/analysis.ws.js} & Connexion WebSocket au serveur d'analyse vidéo. \\
\texttt{store/auth.store.js} & Stockage de l'identité utilisateur, des droits, de l'access token et du refresh token. \\
\texttt{router/index.js} & Routes frontend et gardes de navigation par rôle. \\
\bottomrule
\end{longtable}

\subsubsection{Gestion des rôles côté interface}

Trois profils sont représentés :

\begin{itemize}
  \item \textbf{sportif} : consultation du tableau de bord, lancement de séance, historique, comparaison de deux séances, analyse vidéo ;
  \item \textbf{coach} : consultation des sportifs suivis, accès au live d'une séance active, comparaison de séances, accès aux analyses du périmètre affecté ;
  \item \textbf{administrateur} : gestion des utilisateurs, modules et affectations coach/sportif.
\end{itemize}

Les gardes Vue Router empêchent l'accès aux pages non compatibles avec le rôle local. La sécurité effective reste appliquée côté API.

\subsubsection{Gestion des tokens}

Le frontend conserve deux jetons :

\begin{itemize}
  \item un access token JWT court, signé par l'API et envoyé en priorité dans l'en-tête \texttt{Authorization: Bearer} ;
  \item une copie de compatibilité envoyée dans \texttt{x-session-id} ;
  \item un refresh token long, opaque, utilisé uniquement sur \texttt{POST /auth/refresh}.
\end{itemize}

Lorsqu'une requête API renvoie \texttt{401}, le client Axios tente un refresh. Si le refresh réussit, l'access token est remplacé et la requête initiale est rejouée. Si le refresh échoue, les données locales d'authentification sont supprimées et l'utilisateur est redirigé vers la page de connexion.

\subsection{API Node.js / Express}

\subsubsection{Organisation}

L'API est organisée selon une structure classique routes/contrôleurs/modèles :

\begin{longtable}{p{5cm}p{9cm}}
\toprule
\textbf{Répertoire} & \textbf{Rôle} \\
\midrule
\texttt{routes/} & Déclaration des routes HTTP sous le préfixe \texttt{/trackrapi}. \\
\texttt{controllers/} & Logique métier : authentification, utilisateurs, modules, mesures, séances, IA. \\
\texttt{models/} & Modèles Mongoose correspondant aux collections MongoDB. \\
\texttt{utils/} & Calculs statistiques : distance GPS, pas, stress, score de performance. \\
\texttt{ai/} & Scripts Python d'entraînement et de prédiction. \\
\texttt{services/} & Appels TCP vers le serveur central, exécution des scripts IA, surveillance de sessions. \\
\bottomrule
\end{longtable}

\subsubsection{Format de réponse}

Les réponses API suivent le format :

\begin{lstlisting}[style=trackrcode]
{
  "error": 0,
  "status": 200,
  "data": { ... }
}
\end{lstlisting}

Une valeur \texttt{error} différente de zéro indique une erreur applicative. Le champ \texttt{status} reprend le statut HTTP attendu et \texttt{data} contient les données ou le message d'erreur.

\subsubsection{Authentification et sécurité}

L'authentification repose sur un couple access token JWT / refresh token :

\begin{itemize}
  \item \texttt{POST /auth/signin} vérifie le mot de passe avec bcrypt ;
  \item l'API crée un identifiant de session \texttt{jti}, stocké en base dans le champ \texttt{sessionId} ;
  \item l'access token est un JWT signé avec \texttt{JWT\_SECRET}. Il contient notamment \texttt{sub}, \texttt{login}, \texttt{rights}, \texttt{jti} et \texttt{exp} ;
  \item \texttt{verifyToken} vérifie la signature JWT puis contrôle que le couple \texttt{sub}/\texttt{jti} existe toujours en base ;
  \item le refresh token reste opaque et stocké en base dans le champ \texttt{refreshToken} ;
  \item \texttt{POST /auth/refresh} renouvelle le JWT, le \texttt{jti} et le refresh token, ce qui invalide les anciens jetons ;
  \item \texttt{POST /auth/logout} supprime le \texttt{jti} et le refresh token côté serveur.
\end{itemize}

Les routes internes appelées par le serveur central sont protégées par l'en-tête \texttt{x-service-secret}. Cette protection évite qu'une écriture de mesure ou un auto-enregistrement de module soit accepté depuis un client HTTP externe non autorisé.

\subsubsection{Routes principales}

\begin{longtable}{p{2cm}p{4.8cm}p{3cm}p{5cm}}
\toprule
\textbf{Méthode} & \textbf{Chemin} & \textbf{Protection} & \textbf{Rôle} \\
\midrule
GET & \texttt{/health} & publique & Vérification de disponibilité de l'API. \\
POST & \texttt{/auth/signin} & publique & Connexion utilisateur. \\
POST & \texttt{/auth/refresh} & refresh token & Renouvellement des jetons. \\
POST & \texttt{/auth/logout} & access token & Déconnexion serveur. \\
GET & \texttt{/auth/me} & access token & Lecture de la session courante. \\
POST & \texttt{/measure/create} & service secret & Création de mesure depuis le serveur central. \\
PATCH & \texttt{/measure/update} & admin & Mise à jour d'une mesure. \\
GET & \texttt{/measure/get} & access token & Lecture filtrée des mesures. \\
POST & \texttt{/module/register} & service secret & Auto-enregistrement d'un module matériel. \\
POST & \texttt{/module/connection} & service secret & Mise à jour de l'état de connexion module. \\
POST & \texttt{/module/create} & admin & Création administrative d'un module. \\
PATCH & \texttt{/module/update} & admin & Modification administrative d'un module. \\
GET & \texttt{/module/get} & access token & Liste des modules. \\
POST & \texttt{/user/create} & admin & Création utilisateur. \\
PATCH & \texttt{/user/update} & admin & Modification utilisateur et affectation coach. \\
GET & \texttt{/user/getusers} & admin/coach & Liste des utilisateurs selon périmètre. \\
POST & \texttt{/session/start} & access token & Démarrage d'une séance. \\
POST & \texttt{/session/stop} & access token & Arrêt d'une séance et calcul des statistiques. \\
POST & \texttt{/session/active} & service secret & Vérification interne d'une session active. \\
POST & \texttt{/session/active-for-module} & access token & Vérification frontend d'un module actif. \\
GET & \texttt{/session/history} & access token & Historique filtré par rôle. \\
GET & \texttt{/analysis} & access token & Historique filtré des analyses vidéo. \\
GET & \texttt{/analysis/:analysisId} & access token & Détail d'une analyse vidéo. \\
POST & \texttt{/ai/train} & admin & Ré-entraînement du modèle IA. \\
GET & \texttt{/ai/predict/:sessionId} & access token & Prédiction IA pour une séance accessible. \\
GET & \texttt{/ai/insights/:sessionId} & access token & Explications IA d'une séance accessible. \\
\bottomrule
\end{longtable}

\subsubsection{Documentation OpenAPI / Swagger}

Une documentation Swagger est servie directement par l'API via \texttt{/trackrapi/api-docs}, à partir d'une configuration \texttt{swagger-jsdoc} située dans \texttt{TrackrAPI/config/swaggerConfig.js}. Elle permet de tester les routes publiques, les routes protégées par \texttt{Authorization: Bearer <JWT>}, le renouvellement par refresh token et les routes internes protégées par \texttt{x-service-secret}. Cette documentation sert également de support technique pour expliciter les contrats HTTP du POC.
\subsection{Base de données MongoDB}

\subsubsection{Collections}

\begin{longtable}{p{3.2cm}p{10.8cm}}
\toprule
\textbf{Collection} & \textbf{Description} \\
\midrule
\texttt{users} & Utilisateurs de l'application. Contient login, mot de passe hashé, email, droits, coach référent, access token, refresh token et dates d'expiration. \\
\texttt{modules} & Modules matériels enregistrés. Contient nom, surnom, clé unique, type de microcontrôleur, chipsets, état connecté et date de dernière activité. \\
\texttt{chipsets} & Déclaration des capteurs ou composants associés aux modules. \\
\texttt{sessions} & Séances sportives. Lie un utilisateur, un module, les dates de début/fin, la dernière mesure reçue et les statistiques calculées. \\
\texttt{measures} & Mesures unitaires. Contient type, date, valeur textuelle, module, séance et éventuellement identifiant d'analyse vidéo. \\
\bottomrule
\end{longtable}

\subsubsection{Structure des documents}

Les modèles Mongoose du répertoire \texttt{TrackrAPI/models} décrivent les documents stockés. Les champs principaux sont :

\begin{longtable}{p{3.2cm}p{10.8cm}}
\toprule
\textbf{Collection} & \textbf{Champs principaux} \\
\midrule
\texttt{users} & \texttt{login}, \texttt{password} hashé, \texttt{email}, \texttt{rights}, \texttt{coach}, \texttt{sessionId}, \texttt{sessionExpiresAt}, \texttt{refreshToken}, \texttt{refreshExpiresAt}. Le champ \texttt{coach} référence un utilisateur possédant le rôle coach. \\
\texttt{modules} & \texttt{name}, \texttt{shortName}, \texttt{key}, \texttt{uc}, \texttt{chipsets}, \texttt{connected}, \texttt{lastSeen}. La clé \texttt{key} identifie le module côté TCP. \\
\texttt{chipsets} & \texttt{name}, \texttt{description}, \texttt{links}, \texttt{caps}. Les capacités décrivent les types de mesures ou fonctions disponibles. \\
\texttt{sessions} & \texttt{sessionId}, \texttt{user}, \texttt{module}, \texttt{startDate}, \texttt{lastMeasureAt}, \texttt{lastHeartRateAt}, \texttt{endDate}, \texttt{stats}. Les statistiques agrégées sont stockées dans \texttt{stats}. \\
\texttt{measures} & \texttt{type}, \texttt{date}, \texttt{value}, \texttt{module}, \texttt{session}, \texttt{analysisId}. La valeur reste une chaîne afin d'accepter des mesures hétérogènes. \\
\bottomrule
\end{longtable}

\subsubsection{Relations}

\begin{itemize}
  \item Une séance référence un utilisateur et un module.
  \item Une mesure référence une séance et éventuellement un module.
  \item Un module référence plusieurs chipsets.
  \item Un sportif référence au plus un coach.
  \item Un coach peut être référencé par plusieurs sportifs.
\end{itemize}

\subsection{Serveur central Java}

\subsubsection{Rôle}

Le serveur central Java est le point d'entrée des modules matériels. Il écoute des connexions TCP persistantes, gère l'identification des modules et transmet les mesures vers l'API. Il permet également à l'API de demander le démarrage ou l'arrêt d'une séance sur un module connecté.

\subsubsection{Organisation}

\begin{longtable}{p{5cm}p{9cm}}
\toprule
\textbf{Fichier} & \textbf{Rôle} \\
\midrule
\texttt{DataDriver.java} & Interface décrivant les opérations de stockage, enregistrement module, état session et connexion. \\
\texttt{HttpDataDriver.java} & Implémentation qui transmet les opérations à l'API HTTP. Ajoute le secret de service dans l'en-tête \texttt{x-service-secret}. \\
\texttt{MongoDataDriver.java} & Implémentation historique d'accès direct à MongoDB. \\
\texttt{ThreadServer.java} & Gestion d'une connexion TCP module, parsing des commandes et état local de session. \\
\texttt{MainServer.java} & Démarrage du serveur TCP. \\
\bottomrule
\end{longtable}

\subsubsection{Protocole TCP}

Les commandes principales sont :

\begin{longtable}{p{4cm}p{10cm}}
\toprule
\textbf{Commande} & \textbf{Description} \\
\midrule
\texttt{AUTOREGISTER uc chipsets...} & Demande d'enregistrement automatique d'un module. \\
\texttt{HELLO moduleKey} & Identification d'un module déjà enregistré. \\
\texttt{STOREMEASURE type date value moduleKey} & Envoi d'une mesure pendant une séance active. \\
\texttt{STOREANALYSIS type date value} & Envoi d'un résultat d'analyse vidéo sérialisé. \\
\texttt{START\_SESSION sessionId} & Commande envoyée au module pour démarrer l'acquisition. \\
\texttt{STOP\_SESSION} & Commande envoyée au module pour arrêter l'acquisition. \\
\bottomrule
\end{longtable}

\subsection{Module ESP32}

\subsubsection{Matériel}

Le module matériel est construit autour d'un ESP32. Il utilise :

\begin{itemize}
  \item un GPS connecté sur l'UART matériel 2, avec réception sur la broche GPIO 27 ;
  \item un MPU6050 connecté en I2C, avec SDA sur GPIO 19 et SCL sur GPIO 22 ;
  \item une ceinture cardiaque Bluetooth Low Energy exposant le service Heart Rate \texttt{180D} et la caractéristique \texttt{2A37} ;
  \item le bouton BOOT de l'ESP32 pour une réinitialisation de configuration après maintien prolongé.
\end{itemize}

\subsubsection{Bibliothèques embarquées}

Le firmware PlatformIO utilise notamment :

\begin{itemize}
  \item \texttt{WiFiManager} pour la configuration WiFi et l'adresse du serveur ;
  \item \texttt{MPU6050} pour l'accéléromètre et le gyroscope ;
  \item \texttt{NeoGPS} pour la lecture GPS ;
  \item les bibliothèques BLE de l'ESP32 pour la ceinture cardiaque ;
  \item \texttt{Preferences} pour stocker l'adresse serveur et la clé module en mémoire flash.
\end{itemize}

\subsubsection{Fonctionnement}

Au démarrage, le module charge la configuration stockée. Si aucune configuration WiFi ou serveur n'existe, un point d'accès \texttt{Trackr-ESP32} est ouvert par WiFiManager afin de saisir l'adresse et le port du serveur central. La clé du module est conservée en mémoire flash après auto-enregistrement.

Une fois connecté au serveur central, le module envoie \texttt{AUTOREGISTER} si nécessaire, puis \texttt{HELLO}. Les mesures ne sont envoyées que lorsque le serveur central transmet une commande de démarrage de séance. Les données GPS sont filtrées côté module avec une distance minimale, une vitesse minimale et un intervalle minimal entre deux positions. Le MPU6050 est d'abord stabilisé, puis l'ESP32 échantillonne l'accélération localement afin de détecter les pics de marche/course et maintenir un compteur cumulatif de pas. Ce compteur est transmis comme mesure dédiée \texttt{steps}. Les axes d'accélération et de gyroscope restent envoyés environ une fois par seconde pour le suivi brut et les diagnostics.

\subsection{Calcul des statistiques de séance}

Les statistiques de séance sont calculées à l'arrêt de la séance à partir des mesures enregistrées. Les fonctions principales sont dans \texttt{TrackrAPI/utils}.

\begin{longtable}{p{4cm}p{10cm}}
\toprule
\textbf{Fichier} & \textbf{Rôle} \\
\midrule
\texttt{geo.js} & Calcul de distance par Haversine, filtrage des points GPS aberrants et indicateur de qualité GPS. \\
\texttt{physio.js} & Calcul de stress et estimation de secours des pas si aucune mesure embarquée \texttt{steps} n'est disponible. \\
\texttt{sessionStats.js} & Agrégation globale d'une séance : distance, vitesse, fréquence cardiaque, dernière valeur \texttt{steps}, stress, score et qualité GPS. \\
\texttt{performanceScore.js} & Calcul du score de performance de référence. \\
\bottomrule
\end{longtable}

La distance et la trace GPS restent dépendantes du matériel, de l'environnement et du placement du module. Le comptage de pas est désormais effectué côté ESP32 comme compteur cumulatif basé sur l'accéléromètre, puis stocké sous forme de mesure \texttt{steps}. Lorsque cette mesure existe, les vues utilisent la dernière valeur transmise par l'ESP32 ; l'estimation logicielle depuis les axes \texttt{acc\_x}, \texttt{acc\_y} et \texttt{acc\_z} n'est conservée qu'en solution de secours.

\subsection{Serveur d'analyse vidéo Python}

\subsubsection{Flux d'analyse}

Le serveur d'analyse reçoit la vidéo par WebSocket. Le protocole est constitué de trois messages principaux :

\begin{enumerate}
  \item \texttt{START\_ANALYSIS} : initialise l'analyse, précise l'exercice et l'utilisateur ;
  \item \texttt{VIDEO\_CHUNK} : transmet la vidéo encodée en base64 ;
  \item \texttt{END\_ANALYSIS} : finalise le fichier temporaire et lance l'analyse.
\end{enumerate}

Après analyse, le résultat est renvoyé au client. Le frontend le sauvegarde ensuite via une route API authentifiée afin d'alimenter immédiatement l'historique visible par l'utilisateur. Le serveur d'analyse conserve également le chemin historique de transmission vers le serveur central avec \texttt{STOREANALYSIS}, ce qui permet de rester compatible avec l'architecture initiale.

\subsubsection{Organisation}

\begin{longtable}{p{5cm}p{9cm}}
\toprule
\textbf{Fichier} & \textbf{Rôle} \\
\midrule
\texttt{ws\_server.py} & Serveur WebSocket, réception des messages, orchestration de l'analyse. \\
\texttt{pose\_extractor.py} & Extraction des landmarks MediaPipe sur les frames échantillonnées. \\
\texttt{metrics.py} & Fonctions géométriques : conversion en coordonnées, calcul d'angle, lissage. \\
\texttt{analyzer\_registry.py} & Sélection de l'analyseur selon l'exercice demandé. \\
\texttt{squat\_analyzer.py} & Analyse principale du squat. \\
\texttt{bench\_analyzer.py} & Analyse expérimentale du développé couché. \\
\texttt{deadlift\_analyzer.py} & Analyse expérimentale du soulevé de terre. \\
\texttt{central\_tcp\_client.py} & Envoi du résultat d'analyse au serveur central. \\
\bottomrule
\end{longtable}

\subsubsection{Métriques vidéo}

Le squat est évalué à partir des angles des genoux, de l'inclinaison du buste et de la différence gauche/droite. Les répétitions sont détectées par transition entre une position haute et une position basse. Le développé couché expérimental s'appuie sur les angles des coudes. Le soulevé de terre expérimental s'appuie sur l'extension de hanche.

Les modes expérimentaux dépendent fortement de l'angle de caméra, de la stabilité de l'image et de la visibilité des articulations. Ils constituent une extension fonctionnelle du système, mais nécessitent une validation sur davantage de vidéos avant un usage fiable.

\subsection{Intelligence artificielle}

\subsubsection{Objectif}

Le module IA ne produit plus un score abstrait sur 100. Il prévoit une grandeur directement observable : la fréquence cardiaque moyenne, en battements par minute, de la prochaine séance de charge comparable. Cette prévision peut être confrontée à la fréquence cardiaque réellement mesurée lorsque cette séance est réalisée. L'application conserve alors la valeur prévue, la valeur observée et l'erreur absolue en bpm. Le modèle devient ainsi évaluable et peut progressivement intégrer de nouvelles paires issues du sportif.

La prévision alimente une aide à la décision prudente. Une hausse légère de charge n'est proposée que si les indicateurs de récupération sont favorables et si le modèle bat une baseline simple sur le jeu de validation. Dans le cas contraire, l'interface indique explicitement que la prévision reste informative. Il ne s'agit ni d'un avis médical ni d'un générateur autonome de programmes.

\subsubsection{Modèle et cible}

Le modèle utilisé est \texttt{XGBRegressor}. Ce choix convient aux données tabulaires numériques et permet une inférence rapide. Les caractéristiques sont :

\begin{itemize}
  \item \texttt{distanceKm} et \texttt{durationMin} ;
  \item \texttt{hrAvg} et \texttt{hrMax} ;
  \item \texttt{stress} ;
  \item \texttt{steps} et \texttt{paceStepsPerMin} ;
  \item les écarts de distance, pas, stress et fréquence cardiaque par rapport à l'historique ;
  \item le nombre de séances antérieures et le délai depuis la séance précédente.
\end{itemize}

La cible \texttt{next\_comparable\_session\_hr\_avg} est la fréquence cardiaque moyenne effectivement observée lors de la prochaine séance comparable. Pour les données TrackrAI, deux séances sont comparables si elles durent au moins cinq minutes, contiennent un cardio valide, disposent d'une charge mesurable par distance ou pas, et présentent des rapports de durée et de charge raisonnables. Aucune pondération métier ne fabrique la cible.

\subsubsection{Données externes et entraînement}

Le script \texttt{prepare\_pamap2.py} télécharge l'archive officielle PAMAP2 de l'UCI, y compris son archive ZIP imbriquée, puis lit les fichiers Protocol des neuf sujets. Les signaux cardiaques et inertiels sont agrégés en fenêtres consécutives de 60 secondes appartenant à la même activité. Pour chaque paire, \texttt{nextHrAvg} contient la fréquence cardiaque réellement mesurée dans la fenêtre suivante.

Le fichier généré contient 194 paires brutes. Après exclusion des activités sans charge locomotrice exploitable, le dernier entraînement utilise 81 paires PAMAP2 et 24 paires TrackrAI, soit 105 observations réelles. La séparation entraînement/validation est effectuée par sujet avant l'augmentation synthétique. Le bruit faible n'est appliqué qu'au sous-ensemble d'entraînement : aucune copie d'une observation de validation ne peut donc se retrouver dans l'entraînement.

\begin{lstlisting}[style=trackrcode]
target: next_comparable_session_hr_avg
rawSamples: 105
localSamples: 24
externalSamples: 81
validationSamples: 22
maeBpm: 6.172
rmseBpm: 7.907
r2: 0.793
baselineMaeBpm: 5.434
beatsBaseline: false
\end{lstlisting}

La baseline prévoit simplement que la prochaine fréquence cardiaque moyenne sera identique à l'actuelle. Sur la séparation actuelle, sa MAE de 5,434 bpm reste meilleure que celle de XGBoost, égale à 6,172 bpm. Ce résultat est conservé et affiché honnêtement : le modèle n'est pas encore autorisé à piloter automatiquement une augmentation de charge. Le pipeline est néanmoins fonctionnel, reproductible et mesurable ; de nouvelles séances réelles permettront de vérifier si la personnalisation locale améliore progressivement ce résultat.

\subsubsection{Prédiction et vérification}

Le script \texttt{predict\_session.py} refuse les séances de moins de cinq minutes, sans fréquence cardiaque ou sans distance ni pas exploitables. Pour une séance éligible, il renvoie :

\begin{itemize}
  \item la fréquence cardiaque moyenne prévue en bpm ;
  \item l'écart prévu par rapport à la séance actuelle ;
  \item une plage indicative construite à partir de la MAE de validation ;
  \item les caractéristiques utilisées et les métriques du modèle.
\end{itemize}

Lorsqu'une séance comparable suivante se termine, l'API rattache sa fréquence cardiaque observée à la prévision précédente et calcule l'erreur absolue. L'endpoint \texttt{/ai/insights/:sessionId} expose ces informations au sportif, à son coach affecté et à l'administrateur. Les objectifs de durée, distance ou pas restent issus de règles prudentes et transparentes ; ils ne sont augmentés que si la validation du modèle et les indicateurs physiologiques le permettent.

\section{Sécurité applicative}

\subsection{Contrôle d'accès}

Les accès sont filtrés à deux niveaux. Le frontend masque les pages non compatibles avec le rôle utilisateur, puis l'API vérifie systématiquement les droits avant d'accéder aux données sensibles.

Les règles principales sont :

\begin{itemize}
  \item un sportif consulte uniquement ses propres séances, mesures et analyses ;
  \item un coach consulte les données des sportifs qui lui sont affectés ;
  \item un administrateur consulte et modifie l'ensemble des utilisateurs et modules ;
  \item les routes de création de mesures et d'enregistrement module exigent un secret de service ;
  \item l'entraînement IA est réservé à l'administrateur.
\end{itemize}

\subsection{Gestion access token / refresh token}

Le système d'authentification utilise deux jetons. L'access token est un JWT court, signé côté API et transmis dans les requêtes courantes via \texttt{Authorization: Bearer}. Le champ \texttt{jti} du JWT est stocké côté serveur afin de permettre la révocation lors d'un logout ou d'un refresh. Le refresh token est opaque, plus long, stocké en base et permet d'obtenir un nouveau JWT sans nouvelle saisie du mot de passe. Lors d'un refresh, le refresh token et le \texttt{jti} sont également renouvelés, ce qui invalide les anciens jetons.

En environnement de démonstration, les durées sont :

\begin{itemize}
  \item access token : 15 minutes ;
  \item refresh token : 7 jours.
\end{itemize}

\section{Procédures d'exploitation}

\subsection{Installation et lancement}

La procédure recommandée est :

\begin{lstlisting}[style=trackrcode]
git clone <depot-trackrai>
cd TrackrAI
docker compose -f docker-compose.dev.yml up --build
\end{lstlisting}

Le frontend est disponible sur \url{http://localhost:5173}. L'API est disponible sur \url{http://localhost:4567/trackrapi}. Le serveur central TCP est exposé sur le port \texttt{29000}.

\subsection{Arrêt}

\begin{lstlisting}[style=trackrcode]
docker compose -f docker-compose.dev.yml down
\end{lstlisting}

La suppression du volume MongoDB, utile pour repartir d'une base vide, se fait avec :

\begin{lstlisting}[style=trackrcode]
docker compose -f docker-compose.dev.yml down -v
\end{lstlisting}

\subsection{Données de démonstration}

La base initialisée par l'API contient plusieurs comptes de test permettant de présenter les rôles sans manipulation manuelle préalable. Pour la démonstration finale, le compte \texttt{floryan/floryan} est un coach et le compte \texttt{corentin/corentin} est un sportif affecté à ce coach. Des séances historiques sont générées pour Corentin afin de montrer l'historique, le tableau de bord coach et la comparaison entre séances. Les comptes \texttt{admin/admin}, \texttt{coach/coach}, \texttt{test/azer}, \texttt{thomas/thomas} et \texttt{lea/lea} restent également disponibles.

\subsection{Intégration continue GitLab / Jenkins}

Le versioning est réalisé sur le serveur GitLab du département. L'intégration continue est prévue sur Jenkins à partir du \texttt{Jenkinsfile} présent à la racine du dépôt. Le dépôt contient également un fichier \texttt{.gitlab-ci.yml}, utile si un runner GitLab Docker est disponible. Les pipelines s'appuient sur \texttt{docker-compose.ci.yml} afin de valider la configuration Docker, construire les images applicatives et lancer les vérifications automatisées dans un environnement reproductible. La documentation complémentaire est conservée dans \texttt{docs/CI\_GITLAB\_JENKINS.md}. Dans un déploiement réel, les secrets utilisés par la CI doivent être déplacés dans les variables protégées GitLab ou Jenkins.

\subsection{Tests}

Les tests manuels prioritaires sont :

\begin{enumerate}
  \item connexion administrateur et affectation d'un coach à un sportif ;
  \item connexion coach et vérification du filtre \og mes sportifs \fg{}, de l'accès live à une séance active et de la comparaison de séances ;
  \item connexion sportif et vérification du coach référent ;
  \item démarrage et arrêt d'une séance avec un module ESP32 ;
  \item upload d'une vidéo, consultation du résultat et de l'historique intégré sur la même page ;
  \item vérification de la trace GPS et du compteur de pas \texttt{steps} envoyé par l'ESP32 sur une nouvelle séance.
\end{enumerate}

Les tests d'acceptation peuvent être lancés par :

\begin{lstlisting}[style=trackrcode]
cd TrackrAIAcceptanceTests
mvn test
\end{lstlisting}

\section{Limites et perspectives}

\subsection{Capteurs}

Le GPS reste sensible à l'environnement, en particulier en intérieur ou au démarrage du fix. Le comptage de pas repose uniquement sur l'accéléromètre du MPU6050 et dépend donc du placement du boîtier et du mouvement réalisé. Le compteur embarqué rend la valeur cohérente et cumulative pendant une séance, mais il ne prétend pas atteindre la précision d'un podomètre industriel ou d'une montre sportive calibrée.

\subsection{Analyse vidéo}

Le squat est le mode le plus fiable. Les modes développé couché et soulevé de terre sont fonctionnels mais expérimentaux. Une amélioration future consiste à constituer un jeu de vidéos annotées, à valider les seuils par morphologie et à intégrer plusieurs angles de caméra.

\subsection{Intelligence artificielle}

Le modèle cardiaque est entraîné sur 105 paires observées, dont 81 issues de PAMAP2 et 24 de TrackrAI. Malgré un coefficient $R^2$ de 0,793, sa MAE de 6,172 bpm ne bat pas encore la baseline à 5,434 bpm. L'application expose cette limite et interdit au modèle de justifier seul une hausse de charge. L'objectif prioritaire est de collecter davantage de séances comparables par sportif, de suivre l'erreur prévue/réelle et de réévaluer le modèle avec une validation par individu. Un assistant LLM de génération de programmes reste une perspective distincte et ne doit pas être confondu avec ce modèle prédictif.

\subsection{Sécurité et production}

Les secrets présents dans les fichiers Docker Compose sont adaptés à une démonstration. En production, ils doivent être injectés par un gestionnaire de secrets ou des variables d'environnement non versionnées. L'ajout d'une journalisation centralisée, d'une rotation de secrets et d'une politique CORS plus stricte constitue également une amélioration naturelle.

\section{Annexes}

\subsection{Résumé des ports}

\begin{longtable}{p{4cm}p{3cm}p{7cm}}
\toprule
\textbf{Service} & \textbf{Port dev} & \textbf{Description} \\
\midrule
Frontend & 5173 & Interface Vite. \\
API & 4567 & API HTTP Express. \\
Serveur central & 29000 & Port TCP exposé pour les modules. \\
Analyse vidéo & 6000 & WebSocket d'analyse. \\
MongoDB & 27017 & Base de données. \\
\bottomrule
\end{longtable}

\subsection{Schéma électrique minimal}

\begin{longtable}{p{4cm}p{4cm}p{6cm}}
\toprule
\textbf{Élément} & \textbf{Connexion ESP32} & \textbf{Rôle} \\
\midrule
GPS & RX sur GPIO 27 & Réception des trames NMEA. \\
MPU6050 & SDA GPIO 19, SCL GPIO 22 & Accélération et gyroscope. \\
Ceinture cardio BLE & Bluetooth Low Energy & Fréquence cardiaque et intervalles RR. \\
Bouton BOOT & GPIO 0 & Réinitialisation longue de la configuration. \\
\bottomrule
\end{longtable}

\subsection{Liens utiles}

\begin{itemize}
  \item Documentation Docker : \url{https://docs.docker.com/}
  \item Documentation Vue : \url{https://vuejs.org/}
  \item Documentation Express : \url{https://expressjs.com/}
  \item Documentation MongoDB : \url{https://www.mongodb.com/docs/}
  \item Documentation MediaPipe : \url{https://developers.google.com/mediapipe}
  \item Dataset PAMAP2 UCI : \url{https://archive.ics.uci.edu/dataset/231/pamap2+physical+activity+monitoring}
  \item DOI PAMAP2 : \url{https://doi.org/10.24432/C5NW2H}
  \item Documentation XGBoost : \url{https://xgboost.readthedocs.io/}
  \item Documentation PlatformIO : \url{https://docs.platformio.org/}
\end{itemize}

\end{document}
